% Tipo de documento. En este caso es un art�culo, para folios A4, tama�o de la fuente 11pt y con p�gina separada para el t�tulo
\documentclass[a4paper,12pt,titlepage]{article}

% Carga de paquetes necesarios. OrdenesArticle es un paquete personalizado
\usepackage[spanish]{babel} 
\usepackage[T1]{fontenc}
\usepackage[ansinew]{inputenx} 
\usepackage[spanish,cap,cont,title,fancy]{OrdenesArticle}
\usepackage{setspace}
\usepackage{array}
\usepackage{graphicx}
\usepackage{hyperref}
\usepackage{pifont}
\usepackage{listings}
\usepackage[usenames,dvipsnames]{color}
\usepackage{colortbl}
\usepackage{makeidx}
\hypersetup{bookmarksopen,bookmarksopenlevel=3,linktocpage,colorlinks,urlcolor=blue,citecolor=blue,
						linkcolor=blue,filecolor=blue,pdfnewwindow,
						pdftitle={Auditor�a a DIGSOL}, 
						pdfauthor={Juan Andrada, Jose Domingo L�pez, Antonio Mart�n Menor, Francisco Jos� Oteo}}


% Macro para definir una lista personalizada 
\newenvironment{milista}%
{\begin{list}{\textbullet}%
{\settowidth{\labelwidth}{\textbullet} \setlength{\leftmargin}{\dimexpr\labelsep+\labelwidth+5pt}
\setlength{\itemsep}{\dimexpr 0.5ex plus 0.25ex minus 0.25ex}
\setlength{\parsep}{\itemsep}
\setlength{\partopsep}{\itemsep}
\addtolength{\topsep}{-7.5pt}
}}%
{\end{list}}

\begin{document}

% En las p�ginas de portada e �ndices, no hay encabezado ni pie de p�gina
\pagestyle{fancy} 

% Se incluye la portada
\begin{titlepage}
	\begin{center}
  	{\LARGE UNIVERSIDAD DE CASTILLA-LA MANCHA} \\
  	\bigskip
  	{\Large ESCUELA SUPERIOR DE INFORM�TICA} \\
  	\vspace{20mm}
  	\includegraphics[scale=0.45, keepaspectratio]{esi_bw.png} \\
  	\vspace{20mm}
  	{\Huge \textbf{Auditor�a y Seguridad de la Informaci�n}} \\
  	\vspace{10mm}
  	{\LARGE \textsc{\textbf{Auditor�a a DIGSOL}}} \\
  	\vspace{30mm}
  	\today
  	\vspace{30mm}
  	\begin{flushleft}
  		{\large Juan Andrada Romero}\\
  		\vspace{1mm}
  		{\large Jose Domingo L�pez L�pez}\\
  		\vspace{1mm}
  		{\large Antonio Mart�n Menor de Santos}\\
  		\vspace{1mm}
  		{\large Francisco Jos� Oteo Fern�ndez}\\  		
  	\end{flushleft}
	\end{center}
\end{titlepage}

% Se ajusta la separaci�n entre p�rrafos
\parskip=10pt

% Aqui se incluyen los archivos .tex que forman el documento
% -*- coding: utf-8 -*-

\subsection{Problemas relacionados con la Planificaci�n Estrat�gica de TI}

\begin{spacing}{1.4}

% Problemas relacionados con el objetivo de control PO1.1, PO1.2, PO1.4
Consultando al Director General sobre la existencia de un plan estrat�gico de TI, se ha observado que no existe dicho plan. Por tanto, la empresa no est� alineada con las tecnolog�as de la informaci�n, lo que puede suponer, a largo o medio plazo, que la empresa comience a perder presencia en el mercado, ya que para 
seguir siendo competitiva, es necesario investigar nuevas tecnolog�as, actualizar las que ya tiene y aprovechar las oportunidades que ofrece las tecnolog�as de la informaci�n.

Derivado de la inexistencia de un plan estrat�gico de TI que alinee las tecnolog�as con el objetivo de negocio, aparece otro problema, que es la falta de un Departamento de Sistemas de Informaci�n en la organizaci�n de la empresa. Esto provoca que, cuando cada departamento solicita nuevas tecnolog�as, la direcci�n lo aprueba, ya que no hay ning�n especialista que pueda administrar el valor de las TI actuales. Por esta raz�n, la inversi�n en tecnolog�a que realiza la empresa es m�s elevada de lo necesario, lo que puede traducirse en p�rdidas econ�micas importantes a largo plazo. 

% Problemas relacionados con el objetivo de control PO1.6
Por otra parte, como no existe un plan que coordine todo lo relacionado con las tecnolog�as de la informaci�n, no existen portafolios de proyectos que asigne correctamente las prioridades a los diferentes proyectos que maneja la empresa. As�, los desarrolladores inform�ticos (que est�n asociados a los distintos departamentos de la empresa) se limitan a desarrollar las aplicaciones que dichos departamentos necesiten para llevar a cabo su tarea, en el orden en el que se van solicitando. Esto provocar� que no se atiendan a tiempo los proyectos m�s prioritarios para mantener el objetivo de negocio de la empresa y se produzcan p�rdidas econ�micas, adem�s de perder competitividad en el mercado.

\end{spacing}





\end{document}
